\scanpage{57}
\mspage{16}
Donc la donnée de $\Delta$ est un peu plus fine que celle du syst. de racines associé de type $E_6$ (et a fortiori que la donnée d'une forme \unsure{bilinque|bilinéaire|bil. sym.} \inserted{non dég.} sur un $\mathbb Z$-Module $E$ \ldots). On peut interpréter cette donnée suppl. comme la donnée, \deleted{\unsure{pour un syst}} d'un des deux sommets
\diagram{\begin{tikzcd}[ampersand replacement=\&,column sep=small,row sep=tiny]
\alpha\,\bullet\arrow[r,dash]\&\bullet\arrow[r,dash]\&\bullet\arrow[r,dash]\arrow[d,dash]\&\bullet\arrow[r,dash]\&\bullet\,\beta\\
\&\&\bullet
\end{tikzcd}}
$(\alpha,\beta)$ (extrémités longues) ``\uline{du}'' diagramme de Dynkin de $(E,R)$ (i.e. d'un des syst. simples de racines). Mais on aimerait expliciter, en termes de la géométrie interne à \deleted{\unsure{$R$},} un syst. de racines $(E,R)$ de type $E_6$, comment \deleted{\unsure{exactement à}} \inserted{ce supplément} de structure permet de récupérer le graphe cubique $\Delta$.

Définissons
\diagram{\begin{tikzcd}[ampersand replacement=\&]\Delta\arrow[r,"i"]\&E\end{tikzcd}}
par
\[
i(\delta)=3\delta-\lambda\qquad(\in E\text{ car }\delta.\lambda=1,\ \lambda.\lambda=3)
\]
de sorte que \unsure{pr|pour}
\[
\deleted{\unsure{i\delta|i(\delta)}}\ i(\delta)i(\delta')=9\delta.\delta'-6+3=9\delta\delta'-3
\]
donc
\[
\left\{\begin{aligned}
i(\delta)^2&=-12\\
i(\delta)i(\delta')&=\begin{cases}9&\text{si }\delta,\delta'\text{ liés}\\-3&\text{si }\delta\ne\delta'\text{ non liés}.\end{cases}
\end{aligned}\right.
\]
Donc $i$ est une bijection de $\Delta$ dans \inserted{une partie $\Delta'$ de} l'ens $\unsure{\mathscr E|\mathfrak E}_{(12)}$ des éléments \inserted{$x$} de $E$ satisfaisant
\[
x.x\unsure{{}^{\scriptscriptstyle\leftarrow}}=-12
\]
et sur l'image $\Delta'$ de $\Delta$, \deleted{\unsure{On}} relation ``lié'' s'exprime par
\[
x,x'\text{ liés}\Longleftrightarrow x.x'=9
\]
Bien entendu, $W$ invarie $\Delta'$. \deleted{\unsure{Et}} D'ailleurs $R$ est déduit de $\Delta'$ \deleted{comme} par la Condition que \unsure{!|?}
% archive sheet number: 56
